% Auto-generated by docs/build_integrated_paper.py -- do not edit by hand.
% Standalone scaffolding (preamble, nomenclature, limitations, conclusion,
% \setcounter) stripped for \input into the combined manuscript.
\section{Controller Learning, Hybrid Regularization, and MORL Seed Stability}
\label{sec:results2-control}

Block 1 established a deliberately asymmetric surrogate result: the compact v3 surrogate is the more useful rollout environment for reinforcement learning (RL), whereas the calibrated v3.5 RC--NeuralODE is the more accurate predictive digital twin. Block 2 tests the controller-side consequence of that asymmetry on the live BOPTEST \texttt{bestest\_air} runtime environment. The central question is not whether a controller can be trained on a surrogate, but which functional role each surrogate should play inside the learning loop.

The experiments compare five controller families: pure-v3 thermostatic PPO, direct-v3.5 PPO, hybrid-v3/v3.5 PPO, hierarchical DRL (HDRL), and preference-conditioned MORL. All policies are trained on surrogate backends and then evaluated in closed loop against BOPTEST. One asymmetry must be stated up front: unlike thermostatic PPO and HDRL --- which are evaluated strictly zero-shot after surrogate-only training --- MORL additionally includes a short live-BOPTEST finetuning stage, so MORL is used to probe preference-conditioned robustness and seed stability rather than strict zero-shot surrogate transfer. Two targeted 14-day windows are used for the main thermostatic/HDRL comparison: \texttt{peak\_heat\_window} (January, daily-mean ambient $-24.4\,^\circ$C) and \texttt{typical\_heat\_window} (February, $+2.4\,^\circ$C). MORL and PI reference values additionally use the 12-month yearly evaluation protocol. This difference is intentional: the targeted windows expose controller-family mechanisms, while the yearly protocol exposes seed stability and preference robustness.

Figure~\suppref{S8} summarizes the pipeline. Direct v3.5 is a negative control, not a failed implementation: it asks whether the highest-fidelity predictive surrogate can be used directly as the policy rollout environment (it cannot). The hybrid backend then asks whether the same physical model becomes useful when its role changes from dynamics provider to reward-shaping censor --- yes for thermostatic PPO, no for HDRL at the same $\lambda_T$, and conditionally yes for MORL once the observation interface is widened from 5D to 17D.

\paragraph{Roadmap boundary and executed path.}
Block 2 evaluates which functional role each surrogate should play inside the learning loop, strictly in the order of \texttt{roadmap.md} Block 2:
\begin{enumerate}
  \item train the pure-v3 thermostatic PPO baseline (\S4);
  \item run the direct-v3.5 warm-start negative control (\S4.5);
  \item run the thermostatic hybrid $\lambda_T$ sweep (\S5) and the three-step transfer-gap diagnostics (\S5.5);
  \item run the HDRL $\lambda_T$ sweep (\S6);
  \item run the MORL 5D$\to$17D observation ablation, Pareto sweep, and N=5 canonical seed analysis (\S6.5--\S9);
  \item benchmark the PI reference and rebuild the Block 2 evidence tables (\S10--\S11).
\end{enumerate}
Artifact provenance and rebuild commands for every table and figure are catalogued in roadmap Section 11.1.



This pipeline is only feasible because the surrogate backends are orders of magnitude faster than the live emulator (Fig.~\ref{fig:runtime_feasibility}): millions of policy-gradient steps that would take days against BOPTEST complete in hours in-process.

\begin{figure}[pos={!ht}]
  \centering
  \includegraphics[width=0.70\linewidth]{block2_runtime_fidelity.pdf}
  \caption{\textbf{Surrogates make policy-gradient training computationally feasible.} Environment throughput (env steps/s, log scale) versus 24\,h rollout RMSE$_T$ for the BOPTEST emulator (ground truth) and the three in-process surrogate backends. The surrogates run $85$--$220\times$ faster than the BOPTEST RTE HTTP testbed, turning a days-long PPO run into hours; the hybrid is plotted at v3's rollout RMSE because it rolls out on v3. Data-derived from measured throughput (\texttt{reports/speed\_benchmark\_table.csv}, single benchmark run, CPU, 100 episodes of 96 steps) and the 24\,h rollout-RMSE artefacts.}
  \label{fig:runtime_feasibility}

{\footnotesize\itshape Data: \texttt{reports/\allowbreak speed\_\allowbreak benchmark\_\allowbreak table.csv}}
\end{figure}



Block 2 is structured as a sequence of falsifiable hypotheses, each tied to a roadmap section and resolved against live BOPTEST evidence (Table~\ref{tab:hypotheses}). This pre-specification-style framing is what lets the negative results (direct v3.5, HDRL over-regularization, MORL seed variance) be read as decisions rather than tuning accidents.

\begin{table}[pos={!ht}]
\centering
\caption{Block 2 falsifiable-hypothesis ledger. Each hypothesis is tested on the live BOPTEST RTE in the indicated roadmap section.}
\label{tab:hypotheses}
\small
\begin{tabular}{>{\raggedright\arraybackslash}p{5mm}>{\raggedright\arraybackslash}p{66mm}l>{\raggedright\arraybackslash}p{40mm}}
\toprule
 & Hypothesis & roadmap & Verdict \\
\midrule
H1 & The higher-fidelity twin (v3.5) used directly as the rollout environment yields the best controller. & \S4--4.5 & FALSIFIED (direct v3.5 $m_s>1.0$) \\
H2 & Calibrated v3.5 is useful when its role changes from dynamics provider to a frozen reward-shaping censor. & \S5 & SUPPORTED (hybrid gives best cross-window robustness) \\
H3 & The thermostatic-optimal censor weight $\lambda_T=0.10$ transfers to the HDRL family. & \S6 & FALSIFIED (HDRL best at $\lambda_T=0$) \\
H4 & MORL viability is determined by the reward scalarization alone. & \S6.5--7 & FALSIFIED (needs the 17D observation) \\
H5 & The single-seed MORL canonical ($m_s\approx0.10$) is representative. & \S8--9 & FALSIFIED at N=5 (best of five; CV $0.42$--$0.61$) \\
\bottomrule
\end{tabular}
\end{table}

\subsection{PPO interface, observation design, and reward}

All Block 2 policies use Proximal Policy Optimization (PPO). The shared PPO settings are learning rate $3\times 10^{-4}$ during surrogate pretraining, discount factor $\gamma=0.99$ (an effective horizon of $\sim$25 h at the 900 s step), and ten optimization epochs per rollout. The families differ in rollout length, minibatch size, and total timestep budget because the training scripts were developed per family rather than as a single hyperparameter ablation.



\paragraph{Thermostatic PPO baseline (roadmap \S4).} Block 2 opens with the control baseline: a single-level thermostatic PPO controller trained purely on the v3 rollout surrogate (the canonical 1-hour-step checkpoint \texttt{rc\_node\_v3\_tsupply.pt}) with no v3.5 regularization ($\lambda_T=\lambda_P=0$). The actor and critic are separate MLP heads (Stable-Baselines3 \texttt{MlpPolicy}): the actor maps the 17D observation \eqref{eq:obs17} to a tanh-squashed continuous supply-temperature command \eqref{eq:action_map}, while the critic estimates the state value entering the PPO advantage. Training runs for 10M surrogate steps under the shared settings of Table~\suppref{S13}; with entropy coefficient zero the evaluated policy is deterministic. Crucially, the trained policy is run \emph{zero-shot} on the live BOPTEST RTE with no live finetuning, so the baseline is already a strict surrogate-to-simulator transfer test. It is the reference against which the two negative controls (direct v3.5, warm-start) and the hybrid are measured.

The canonical observation interface is the 17-dimensional extended TSup-style vector (full per-feature breakdown in Supplementary Table~\suppref{S14}):
\begin{equation}
  s_t =
  \left[
  x^{\mathrm{phys}}_t,\,
  x^{\mathrm{time}}_t,\,
  \widehat{T}^{\mathrm{amb}}_{t+1:t+24h},\,
  a_{t-1},\,
  \Delta T^{\mathrm{zone}}_t
  \right] \in \mathbb{R}^{17}.
  \label{eq:obs17}
\end{equation}
It contains 5 physical states (zone temperature, CO$_2$, clipped-log power, previous supply temperature, ambient), 4 cyclic time features, 5 ambient forecasts ($+1,+3,+6,+12,+24$ h), 2 previous-action terms, and 1 causal-smoothed $\Delta T_{\mathrm{zone}}$. The failed MORL baseline uses only the earlier 5D observation, which lacks sufficient actuation and forecast context.



The action is a single normalized supply-temperature command,
\begin{equation}
  a_t \in [-1,1],
  \qquad
  T^{\mathrm{sup}}_t = 18 + \tfrac{a_t+1}{2}(35-18)\quad [^\circ\mathrm{C}],
  \label{eq:action_map}
\end{equation}
with a 1.0 C deadband and a per-step rate limit; the comfort band is $21$--$24\,^\circ$C.

\paragraph{Evaluation metric.} The headline maintenance score combines the comfort-violation rate and the worst-case relative severity over a rollout,
\begin{equation}
\begin{aligned}
  m_s &= r_{\mathrm{time}} + r_{\mathrm{sev}},
  \qquad
  r_{\mathrm{time}} = \frac{1}{N}\sum_{t=1}^{N} \mathbf{1}\!\left[\,T_t < T_{\ell}\ \text{or}\ T_t > T_{h}\,\right], \\[3pt]
  r_{\mathrm{sev}} &= \max_{t}\, \max\!\left(\frac{(T_{\ell}-T_t)_+}{T_{\ell}},\ \frac{(T_t-T_{h})_+}{T_{h}}\right).
\end{aligned}
\label{eq:ms}
\end{equation}
with $T_{\ell}=21\,^\circ$C and $T_{h}=24\,^\circ$C (source: \texttt{evaluation/benchmark\_bestest\_air\_article7\_style.py}). Hence $r_{\mathrm{time}}$ is the fraction of steps outside the band (violation\,$\%=100\,r_{\mathrm{time}}$) and $r_{\mathrm{sev}}$ is the single worst relative band exceedance; lower $m_s$ is better. RMSE$_T$ is reported against the band center $T^{\star}=22.5\,^\circ$C.



The per-step comfort term is piecewise linear with an asymmetric, ambient-dependent slope (source \texttt{envs/backends/boptest\_backend.py}):
\begin{equation}
  r^{\mathrm{comfort}}_t =
  \begin{cases}
    -w_u\,(T_{\ell} - T_t) + b_{\mathrm{heat}} & T_t < T_{\ell}, \\
    +\beta & T_{\ell}+\delta \le T_t \le T_h-\delta, \\
    -w_o\,(T_t - T_h) + b_{\mathrm{cool}} & T_t > T_h,
  \end{cases}
  \label{eq:comfort}
\end{equation}
where $w_u=1.60$ when $T_{\mathrm{amb}}\le 8\,^\circ$C (else $1.15$), $w_o=1.80$ when $T_{\mathrm{amb}}\ge 24\,^\circ$C (else $1.15$), the in-band bonus is $\beta=0.05$, the deadband is $\delta=0.5\,^\circ$C, and $b_{\mathrm{heat}},b_{\mathrm{cool}}$ are small action-direction bonuses. The energy term is $r^{\mathrm{energy}}_t=-\eta P_t$ with $\eta=2\times10^{-4}$, and MORL scalarizes the objectives as $r_t = w_c\,r^{\mathrm{comfort}}_t + w_e\,r^{\mathrm{energy}}_t + w_s\,r^{\mathrm{safety}}_t$. The asymmetric cold/hot weights encode that recovering a cold zone in cold weather is physically harder than the symmetric reverse.

\subsection{Hybrid backend: mathematical role of v3.5}

The hybrid backend changes the role of the calibrated physical surrogate. Instead of rolling out the policy on v3.5 directly, PPO rolls out on v3 and evaluates frozen v3.5 in parallel on the same state-action pair. The per-step reward is augmented by a disagreement penalty:
\begin{align}
  r^{\mathrm{hyb}}_t
  &= r^{\mathrm{comfort}}_t + r^{\mathrm{smooth}}_t + r^{\mathrm{energy}}_t
  - \lambda_T \left|T^{v3}_{t+1}-T^{v3.5}_{t+1}\right|
  - \lambda_P \left|P^{v3}_{t+1}-P^{v3.5}_{t+1}\right|.
  \label{eq:hybrid_reward}
\end{align}
For the canonical thermostatic hybrid, $\lambda_T=0.10$ and $\lambda_P=5.0\times 10^{-5}$; PPO otherwise computes the advantage $A_t = r^{\mathrm{hyb}}_t + \gamma V(s_{t+1}) - V(s_t)$ unchanged. Thus v3.5 is neither a second policy loss nor a direct dynamics model: it is a frozen physics-informed censor that discourages the policy from entering state-action regions where the smooth v3 rollout and the calibrated physical model disagree. Across the canonical hybrid traces (\texttt{reports/hybrid\_disagreement\_summary.csv}, overall), the mean temperature disagreement is $0.969\,^\circ$C (p95 $2.52\,^\circ$C) and the mean power disagreement is $708$ W (p95 $1236$ W) --- large enough to shape learning, bounded enough to stay informative.



\subsection{Thermostatic PPO: direct-v3.5 failure and hybrid success (roadmap \S4--\S5)}

Table~\ref{tab:main_kpi} and Figure~\suppref{S10} summarize the main Block 2 controller result. Pure v3 PPO is already usable ($m_s=0.073$ peak, $0.095$ typical). Direct v3.5 PPO fails catastrophically despite v3.5's superior Block 1 predictive fidelity: live violation reaches 77.1\% on peak and 82.4\% on typical, with RMSE above $4.3\,^\circ$C. The hybrid backend resolves the conflict: $m_s=0.087$ on peak and $0.041$ on typical, with violation below $5\%$ on both windows and lower energy than pure v3 on the peak window.

\begin{table}[pos={!ht}]
\centering
\caption{Canonical Block 2 live BOPTEST controller comparison on the targeted 14-day windows. Data sources for every row are catalogued in roadmap Section 11.1.}
\label{tab:main_kpi}
\small
\begin{tabular}{llrrrr}
\toprule
Policy/backend & Scenario & $m_s$ & Violation (\%) & RMSE$_T$ ($^\circ$C) & Energy (kWh) \\
\midrule
pure v3 PPO & peak & 0.073 & 1.49 & 0.869 & 322.2 \\
direct v3.5 PPO & peak & 1.046 & 77.08 & 4.320 & 556.2 \\
hybrid $\lambda_T=0.10$ & peak & 0.087 & 4.69 & 0.795 & 305.3 \\
pure v3 PPO & typical & 0.095 & 4.39 & 0.622 & 368.3 \\
direct v3.5 PPO & typical & 1.102 & 82.37 & 4.401 & 580.5 \\
hybrid $\lambda_T=0.10$ & typical & 0.041 & 2.38 & 0.633 & 352.8 \\
\bottomrule
\end{tabular}
\end{table}

\begin{figure}[pos={!ht}]
  \centering
  \includegraphics[width=0.78\linewidth]{block2_fidelity_utility_scatter.pdf}
  \caption{The fidelity--utility paradox, separated into the single-model effect and the hybrid resolution. \textbf{(A)} The three single-model surrogates: a \emph{lower} 24\,h rollout RMSE$_T$ (more predictive fidelity, $x$) goes with a \emph{higher} live maintenance score $m_s$ ($y$, worse controller). The least accurate hourly v3 ($1.557\,^\circ$C) is the only usable training environment ($m_s\approx0.08$), while the more accurate matched-resolution v3 ($0.876\,^\circ$C) and calibrated v3.5 ($0.644\,^\circ$C) both collapse above the $m_s=1$ line. \textbf{(B)} The role-separating hybrid breaks the trade-off: it is plotted at v3's rollout RMSE because v3 supplies the transition dynamics while the frozen v3.5 enters \emph{only} as a per-step reward censor, yet it lands in the usable band. Filled markers = peak window, open = typical window (shown separately, not a vertical segment); error bars = $\pm1$ seed s.d.\ over $N=3$ seeds for v3, matched-v3 and the hybrid (\texttt{reports/block2\_thermostatic\_seed\_band.csv}), while direct v3.5 is single-seed. Dashed red line = $m_s=1$ collapse; shaded band = usable region ($m_s<0.1$).}
  \label{fig:paradox_scatter}
\end{figure}



\paragraph{Reading $m_s$ step by step.} The decomposition $m_s=r_{\mathrm{time}}+r_{\mathrm{sev}}$ (Eq.~\ref{eq:ms}) exposes the failure structure directly (Figure~\suppref{S11}). Direct v3.5's peak $m_s=1.046$ is dominated by $r_{\mathrm{time}}\approx 0.77$ (out of band 77\% of the time) plus a large worst-case severity $r_{\mathrm{sev}}\approx 0.28$; the hybrid's peak $m_s=0.087$ is almost entirely a small $r_{\mathrm{sev}}$ with $r_{\mathrm{time}}<0.05$. The hybrid therefore does not merely lower the average error --- it removes the sustained band departures that dominate the direct-v3.5 score.



The time-series and action diagnostics in Figures~\suppref{S12} and~\ref{fig:action_phase} explain the mechanism. Direct v3.5 learns a bang-bang-like control law that drives the live simulator outside the comfort band; in phase space it places extreme actions in temperature-error regimes the live building does not support. The \emph{full training-trajectory loss-landscape} mechanism is not directly measured here; however, the surrogate-side response-surface roughness and the policy-side action saturation that drive this behaviour \emph{are} measured directly in Section~\ref{ssec:coarse-graining} (Table~\ref{tab:surface_sharpness}, Figure~\ref{fig:surface_curves}). What remains hypothesized is which micro-cause dominates (higher advantage-estimator variance under sharper surrogate predictions vs.\ overfitting to sub-step physical structure unusable at the 15-min cadence); discriminating the two is deferred to future work.



\begin{figure}[pos={!ht}]
  \centering
  \includegraphics[width=0.96\linewidth]{block2_q1_polish_phase_density.pdf}
  \caption{Action-density phase portrait of the three thermostatic controllers on the live BOPTEST \texttt{typical\_heat\_window}: empirical hexbin density of the normalised supply-temperature action $a_0$ (one hue per backend) against the thermal error $T_{\mathrm{zone}}-T_{\mathrm{set}}$, on \emph{identical} axes, with the marginal $a_0$ distribution in the right strip of each panel and the measured saturation share ($|a_0|\geq0.9$). The red bands mark the actuator-saturation region. Direct v3.5 collapses onto $a_0=-1$ (100\% saturation, a bang-bang law) whereas the usable hourly v3 and the hybrid modulate ($\approx$19--20\%). Data: closed-loop traces in \texttt{outputs/.../traces/typical\_heat\_window\_thermostatic.csv}.}
  \label{fig:action_phase}

{\footnotesize\itshape Data: \texttt{outputs/\allowbreak block2\_\allowbreak thermostatic\_\allowbreak hybrid\_\allowbreak v3\_\allowbreak v35\_\allowbreak l010/\allowbreak }}
\end{figure}

\subsection{Temporal-coarse-graining ablation: matched-resolution v3 (roadmap \S4.6)}
\label{ssec:coarse-graining}

The direct-v3.5 failure could be dismissed as an artefact of the grey-box / physics-informed model class rather than of fidelity itself. We close this confound with a matched-config control: we retrain the \emph{same black-box} v3 architecture on the matched fifteen-minute corpus---making it a strictly \emph{more accurate} predictor (24\,h rollout RMSE $0.876\,^\circ$C vs $1.557\,^\circ$C for the canonical hourly v3)---and train a PPO controller on it under the \emph{identical} recipe as the canonical pure-v3 baseline, changing only the surrogate's training resolution. The result is unambiguous (Table~\ref{tab:coarse_graining}): the more accurate matched-resolution v3 \emph{collapses} in closed loop ($m_s=1.142$/$1.211$, 85.6/91.4\% violation on peak/typical), on a par with the direct-v3.5 failure and far worse than the hourly-trained controller.

\begin{table}[pos={!ht}]
\centering
\caption{Matched-resolution closed-loop ablation. Increasing v3's temporal resolution improves its predictive RMSE yet destroys its downstream control utility, isolating temporal coarse-graining---not the black-box versus grey-box model class---as the operative variable in this comparison. All $m_s$ and violation cells are live BOPTEST on the two targeted 14-day windows (peak / typical).}
\label{tab:coarse_graining}
\small
\begin{tabular}{llrrlll}
\toprule
Training backend & Architecture & Step (s) & 24\,h RMSE & $m_s$ (pk/typ) & Viol.\ \% (pk/typ) & Verdict \\
\midrule
canonical v3 & black-box & 3600 & 1.557 & 0.073 / 0.095 & 1.5 / 4.4 & useful \\
matched v3 & black-box & 900 & 0.876 & 1.142 / 1.211 & 85.6 / 91.4 & collapse \\
calibrated v3.5 & grey-box RC & 900 & 0.644 & 1.046 / 1.102 & 77.1 / 82.4 & collapse \\
hybrid & v3 + v3.5 censor & 900 & --- & 0.087 / 0.041 & 4.7 / 2.4 & robust \\
\bottomrule
\end{tabular}
\end{table}

This ablation resolves the timestep confound directly: improving v3's temporal resolution and predictive RMSE does not improve downstream control utility---it destroys it. The useful property of the canonical v3 is therefore not its black-box architecture alone, but the optimization-friendly smoothing induced by temporal coarse-graining. Read across the three single-model backends the pattern is near-monotonic in fidelity: the least accurate surrogate (hourly v3) is the only usable training environment, while the more accurate matched v3 and the most accurate calibrated v3.5 both fail. This reframes the hybrid as an explicit \emph{architectural separation of smoothing and fidelity}: v3 supplies the coarse, optimization-friendly rollout dynamics while the frozen v3.5 supplies a physical-plausibility signal as a censor, so policy-gradient training keeps the smoothing it needs while regaining physical grounding. The broader principle is that PPO needs not only accurate predictions but an optimization-friendly training landscape.

\paragraph{Seed robustness.} To check that the single-seed scores are not seed-luck, the three thermostatic backends (pure v3, the matched-resolution v3, and the canonical hybrid) were retrained and re-evaluated on three fixed seeds (\{42, 43, 44\}). Each stays in its qualitative regime on every seed: the usable controllers (pure v3 and hybrid) keep every individual seed below the $5\%$ comfort-violation bar on both windows, while the matched-resolution v3 \emph{collapses on every seed} ($m_s = 1.14$/$1.21$, std $<0.001$ --- the unlearnable environment drives all seeds to the same degenerate policy). The usable-versus-collapse verdicts, including the temporal-coarse-graining failure, are therefore seed-stable (Supplementary Table~\suppref{S16} and Fig.~\suppref{S13}). The single-seed peak advantage of pure v3 over the hybrid (Table~\ref{tab:main_kpi}) does not persist across seeds: over $N{=}3$ the hybrid is at least as good on both windows. The hierarchical and MORL families are treated separately (HDRL single-seed; MORL over $N{=}5$).





\paragraph{Measured mechanism: response-surface smoothness.} The ``optimization-friendly smoothing'' invoked above is not a post-hoc narrative: it is directly measurable on the surrogates themselves, independently of any controller. For each single-model backend we hold a representative state fixed, sweep the supply-temperature action $a_0$ over $[-1,1]$, and read the surrogate's own one-step map $\hat T(a_0)$; from that curve we compute the slope magnitude $\overline{|\partial\hat T/\partial a_0|}$ and, as a step-length-independent measure of non-smoothness, the \emph{relative roughness} (mean curvature normalised by mean slope, both in $^\circ$C per unit action), averaged over a grid of zone temperatures and the two windows' ambients. The scale-free relative roughness orders cleanly with downstream usability (Table~\ref{tab:surface_sharpness}): the only usable training environment, the canonical hourly v3, exposes the \emph{smoothest}, near-monotone action$\to$next-temperature landscape (relative roughness $0.17$), whereas both backends that collapse in closed loop expose markedly rougher landscapes---$9.4\times$ for the matched-resolution v3 and $7.9\times$ for the calibrated v3.5. (The raw per-step slope and range are intentionally \emph{not} used as the mechanism metric: they are confounded by the step length, which is the coarse-graining variable; the normalised roughness is not.) This is the surrogate-side cause whose policy-side symptom is the near-bang-bang action gap of the collapsed controllers (Fig.~\ref{fig:action_phase}): a rougher, more sharply peaked optimisation landscape is precisely what a policy-gradient learner saturates into a non-transferable extremal law. The contrast is visual as well as numerical---the usable hourly v3 has a smooth, monotone action map, while both collapsing surrogates are visibly non-monotone (Fig.~\ref{fig:surface_curves}), so this figure (cause) pairs directly with the bang-bang phase portrait of Fig.~\ref{fig:action_phase} (effect).

\begin{table}[pos={!ht}]
\centering
\caption{Measured response-surface smoothness of the three single-model backends, probed directly on each surrogate's one-step action map (no controller, no BOPTEST; data: \texttt{reports/block2\_mechanism\_surface\_sharpness.csv}). The scale-free relative roughness (curvature normalised by slope) is independent of the step length and tracks closed-loop usability: the usable hourly v3 is the smoothest; both collapsed backends are several times rougher.}
\label{tab:surface_sharpness}
\small
\begin{tabular}{lrrll}
\toprule
Surrogate backend & Slope $\overline{|\partial\hat T/\partial a_0|}$ & Rel.\ roughness & vs.\ v3 & Closed-loop role \\
\midrule
canonical v3 (hourly) & 3.368 & 0.169 & 1.0 (ref) & usable \\
matched v3 (15-min) & 0.272 & 1.578 & 9.4$\times$ & collapse \\
calibrated v3.5 & 0.467 & 1.338 & 7.9$\times$ & collapse \\
\bottomrule
\end{tabular}
\end{table}

\begin{figure}[pos={!ht}]
  \centering
  \includegraphics[width=0.96\linewidth]{block2_surface_response_curves.pdf}
  \caption{Measured roughness mechanism, four panels (probed directly on each surrogate's one-step action map -- no controller, no BOPTEST). \textbf{(A)} Action$\to$next-temperature response $\hat T(a_0)$: thin lines are the per-state curves over the zone-temperature$\times$ambient grid, the thick line is the per-backend median (centred and scaled to a common range so the comparison is of \emph{shape}; line styles aid colour-blind reading) -- the usable hourly v3 is smooth and monotone, both collapsing backends are visibly non-monotone. \textbf{(B)} Scale-free relative roughness (mean curvature $\div$ mean slope) as a \emph{distribution} over the state grid (box + per-state points; $\times$ = canonical aggregate matching Table~\ref{tab:surface_sharpness}): the collapsing backends are $7.9$--$9.4\times$ rougher than the usable v3. \textbf{(C)} Measured policy-side symptom -- fraction of closed-loop BOPTEST steps with a saturated action ($|a_0|>0.9$): hourly v3 modulates ($24\%$), both collapsing backends are near bang-bang ($100\%$). \textbf{(D)} The live BOPTEST outcome (maintenance score $m_s$, peak window) past the $m_s=1$ collapse line. The ordering is \emph{mechanistically consistent} with the paradox -- rougher surface, more saturation, worse live $m_s$ -- but is an association, not a proven causal law (the full training-trajectory loss-landscape is not measured; see Fig.~\ref{fig:action_phase} for the paired policy-side effect). Data: \texttt{reports/block2\_mechanism\_surface\_sharpness.csv} and the closed-loop traces.}
  \label{fig:surface_curves}
\end{figure}

\paragraph{Hybrid $\lambda_T$ sweep and canonical selection (roadmap \S5).} The canonical hybrid operating point was selected by a thermostatic sweep over the temperature-disagreement weight $\lambda_T\in\{0.05,0.10,0.15\}$ at fixed $\lambda_P=5\times10^{-5}$ (Table~\suppref{S17}). The mid setting $\lambda_T=0.10$ (\texttt{hybrid\_l010}) is canonical: per the roadmap selection rule it retains the energy advantage over pure v3 while avoiding the stronger comfort degradation seen at the weaker ($0.05$) and stronger ($0.15$) censor settings. Only the canonical point's live-BOPTEST KPIs are retained as a frozen artifact; the bracketing points served selection only.



\subsection{Warm-start negative control (roadmap \S4.5)}

A second negative control tests whether v3.5 is useful as a policy \emph{initializer} rather than a reward-shaping censor. The benchmark (\texttt{run\_block2.py warmstart}) runs four steps: \emph{(1) pretrain} a thermostatic policy on the calibrated v3.5 surrogate; \emph{(2) scratch finetune} a thermostatic policy from random initialization directly on the live BOPTEST RTE; \emph{(3) warm-start finetune} a second policy on BOPTEST initialized from the v3.5-pretrained checkpoint of step~1; and \emph{(4) evaluate} both finetuned policies on the same BOPTEST windows and compare. The comparison is therefore internally controlled --- scratch and warm-start differ only in initialization. It does not help: warm-started policies are markedly worse than the scratch finetune (Table~\suppref{S18}), raising $m_s$ by roughly two to three times on both windows. The problem is the role assigned to the surrogate during early policy formation, not a lack of fine-tuning; v3.5 belongs in the reward as a censor, not in the weights as an initializer.





\subsection{Transfer-gap diagnostics (roadmap \S5.5)}

This diagnostic (roadmap \S5.5) compares three controllers --- pure v3, hybrid $\lambda_T=0.10$, and direct v3.5 --- in three steps. \emph{Step 1}: a standalone direct-v3.5 thermostatic policy is trained on the calibrated v3.5 environment (the failure control; \texttt{thermostatic-train --variant v35\_direct}), separate from the warm-start checkpoint of the previous section because it answers a different question (zero-shot live transfer rather than warm-start utility). \emph{Step 2}: all three policies are run in closed loop on the live BOPTEST RTE for the two 14-day windows (\texttt{thermostatic-transfer --variant all}). \emph{Step 3}: for each policy we compute four transfer diagnostics (\texttt{thermostatic-diagnose}) --- the temperature tracking gap (live RMSE$_T$), the surrogate-vs-live maintenance-score mismatch
\begin{equation}
\begin{aligned}
  \Delta m_s &= m_s^{\mathrm{surrogate}} - m_s^{\mathrm{BOPTEST}}
  \qquad(\text{negative}\Rightarrow\text{surrogate optimistic}), \\[3pt]
  g_a &= \frac{1}{N}\sum_{t=1}^{N}\big\|\,a_t^{\mathrm{BOPTEST}}-a_t^{\mathrm{surrogate}}\,\big\|_{2} .
\end{aligned}
\label{eq:transfer_gap}
\end{equation}
the action-gap norm $g_a$, and the first divergence step (earliest 15-min step where surrogate and live trajectories separate beyond threshold), with the top driver feature attributed to the divergence.

Direct v3.5 has the largest transfer mismatch ($|\Delta m_s|\approx 0.9$--$1.0$, $g_a\approx 2.0$, live RMSE$_T>4.3\,^\circ$C), and its top divergence driver is $t_{\mathrm{zone}}$: its sharp temperature dynamics produce live actions unlike the surrogate rollout. Hybrid $\lambda_T=0.10$ has the smallest gap ($|\Delta m_s|\approx 0.02$) and, on the typical window, holds the BOPTEST-consistent action for 16 steps before drifting (Table~\suppref{S19}). The action-gap norm is the most diagnostic column: $g_a=2.0$ for direct v3.5 means the policy sits at the opposite action bound from what BOPTEST would choose throughout the episode, i.e.\ a learned bang-bang exploit of the surrogate.





\subsection{HDRL sensitivity: the hybrid weight is controller-family specific (roadmap \S6)}

HDRL here is a \emph{seasonal} hierarchy. A high-level seasonal selector routes control to one of two low-level PPO setpoint specialists --- a winter agent (trained 5M steps on cold-season day ranges with cold-biased comfort shaping) and a summer agent (7M steps on warm-season ranges) --- each acting on the same 17D observation and supply-temperature action. Formally the controller is $\pi(a_t\mid s_t)=\pi_{k(t)}(a_t\mid s_t)$ with the high-level gate $k(t)\in\{\text{winter},\text{summer}\}$ chosen by season; the hierarchy supplies regime-aware temporal abstraction, so each low-level specialist solves a narrower, better-conditioned control problem than a single global policy (Figure~\suppref{S16}).



The HDRL experiment asks whether the thermostatic hybrid setting $\lambda_T=0.10$ transfers to this hierarchical controller. It does not (Table~\suppref{S20}). The mechanism is over-regularization: because each seasonal specialist already encodes comfort-aware structure (season-tuned shaping plus the high-level gate), the additional v3.5 temperature-disagreement censor constrains an already comfort-aware low-level loop, biasing it toward conservative under-heating; performance therefore degrades monotonically with $\lambda_T$, and the correct transfer keeps only the power channel ($\lambda_T=0$). HDRL performs best at $\lambda_T=0.00$ on both windows and degrades monotonically as temperature-disagreement regularization is increased. This shows the correct physical-censor strength depends on the controller family and its action decomposition, not on a universal weight (Fig.~\ref{fig:lambda_specificity}).

\begin{figure}[pos={!ht}]
  \centering
  \includegraphics[width=0.70\linewidth]{block2_lambda_specificity.pdf}
  \caption{\textbf{The optimal censor weight $\lambda_T$ is controller-family specific.} Live BOPTEST maintenance score $m_s$ (cross-window mean of the peak and typical windows) versus $\lambda_T$ for each controller family; $\star$ marks the per-family optimum. Thermostatic PPO improves with the censor (optimum $\lambda_T=0.10$, the canonical hybrid), HDRL is best \emph{without} it ($\lambda_T=0.00$) and degrades monotonically, and MORL adopts $\lambda_T=0.00$. HDRL has a dense four-point sweep; PPO is shown at its two measured endpoints (pure v3 $=\lambda_T0$, hybrid $=0.10$) and MORL at its single adopted configuration. The panel compares the \emph{location} of each optimum, not absolute $m_s$ across families (they use different observation/action interfaces). Data: \texttt{reports/block2\_hdrl\_lambda\_sweep\_summary.csv}, \texttt{block2\_fidelity\_utility\_scatter.csv}, \texttt{block2\_morl\_comparison\_summary.csv}.}
  \label{fig:lambda_specificity}

{\footnotesize\itshape Data: \texttt{reports/\allowbreak block2\_\allowbreak hdrl\_\allowbreak lambda\_\allowbreak sweep\_\allowbreak summary.csv}}
\end{figure}



\paragraph{Engineering reading.} As $\lambda_T$ rises, energy falls slightly (the censor biases the low-level loop toward conservative under-heating, e.g. $329.6\to300.4$ kWh on peak) but comfort collapses --- peak violation roughly quadruples ($6.1\%\to23.0\%$) and RMSE$_T$ nearly doubles. HDRL already supplies comfort-aware temporal abstraction through its scheduler, so an additional temperature-disagreement censor is redundant and over-constrains the setpoint controller; the correct transfer of the Block 1 recipe to HDRL keeps only the power channel ($\lambda_T=0$, $\lambda_P=5\times10^{-5}$).



\subsection{MORL observation ablation: 5D failure to 17D success (roadmap \S6.5--\S7)}

MORL uses a four-stage pipeline that differs from the single-stage PPO families (roadmap Section 7): (1) a 2M-step surrogate \emph{pretrain} on the 17D hybrid backend with canonical $(w_c,w_e,w_s)=(0.80,0.20,0.00)$; (2) an \emph{ERAM} weight-adaptation stage (20 iterations of 100k steps from initial weights $0.34/0.33/0.33$, weight-update temperature $\tau_w=0.35$); (3) a 100k-step \emph{finetune} on the live BOPTEST RTE at learning rate $10^{-4}$ with $\pm3$-day episode-start jitter; and (4) a 12-month \emph{yearly evaluation}. MORL is the only family with a live-BOPTEST finetune; the thermostatic/HDRL families are evaluated zero-shot after surrogate-only training, which is a strictly harder transfer (Figure~\suppref{S18}).



MORL initially failed with a 5D observation (zone temperature, ambient, hour, day, occupancy). Under the current code path, a reconstructed 5D rerun obtains RMSE$_T=2.721\,^\circ$C, violation $37.8\%$, and $m_s=0.680$ (an earlier historical 5D artifact, retained only for audit, is reported in the supplementary). Replacing the observation with the 17D TSup-style vector recovers a usable policy: RMSE$_T=0.720\,^\circ$C, violation $4.9\%$, $m_s=0.099$ (Table~\ref{tab:morl5d17d}). The dominant MORL bottleneck was the observation geometry, not the reward scalarization alone. Widening the observation to include action history and short-horizon forecasts is the standard remedy for the partial observability (POMDP structure) of the underlying control task --- the richer interface restores an approximately Markov state, consistent with forecast-augmented RL-HVAC control~\citep{Gao2024Predictive}.

\begin{table}[pos={!ht}]
\centering
\caption{MORL observation-interface ablation. The backend is hybrid in both cases; only the observation interface changes. The current-code reconstructed 5D rerun is the reproducible main-paper evidence (roadmap Section 11.1).}
\label{tab:morl5d17d}
\small
\begin{tabular}{lrrrr}
\toprule
Variant & Obs dim & RMSE$_T$ ($^\circ$C) & Violation (\%) & $m_s$ \\
\midrule
MORL 5D (current-code rerun) & 5 & 2.721 & 37.8 & 0.680 \\
MORL 17D power-only (canonical) & 17 & 0.720 & 4.9 & 0.099 \\
\bottomrule
\end{tabular}
\end{table}



\paragraph{Power-only backend and claim discipline (roadmap \S7, \S9, \S13).} MORL inherits the controller-family lesson of the HDRL sweep: it runs on the \emph{power-only} hybrid backend ($\lambda_T=0$, $\lambda_P=5\times10^{-5}$), because temperature-disagreement censoring hurts non-thermostatic families. The MORL claim is deliberately narrow and audit-protected. Before the seed expansion, commit \texttt{93df9b3} froze the canonical plan and the seed-45/46 falsification predictions; after $N=5$, commit \texttt{62dc859} appended the observed variance without rewriting the pre-specification, and a replay test confirmed bit-identical BOPTEST trajectories for a fixed checkpoint (so the spread is training stochasticity, not simulator noise). The defensible reading is therefore: under the fixed 17D power-only backend, MORL is substantially stronger than the reconstructed 5D interface and yields a usable comfort--energy Pareto structure, but the $N=5$ canonical variance (CV $0.42$--$0.61$) is too high for a deployment-stability claim --- the single-seed canonical is the best of five, not the median.

\subsection{MORL Pareto front and N=5 seed variance (roadmap \S8--\S9)}

The MORL Pareto sweep varies comfort/energy weights with safety weight zero (Table~\suppref{S21}). Energy-only control collapses comfort ($m_s=1.588$, violation $86.8\%$); comfort-leaning settings are far more stable. The single-seed (seed 42) Pareto points are reported separately from the $N=5$ canonical extensions, because the seed analysis is the central audit result: the neutral 50/50 canonical has $m_s=0.187\pm0.078$ (CV $0.42$, 95\% $t$-CI $[0.090,0.284]$) and the practical 75/25 canonical has $m_s=0.139\pm0.085$ (CV $0.61$, 95\% $t$-CI $[0.033,0.245]$). Seed 46 is an outlier in both groups (Table~\suppref{S22}). Because the replay audit produced bit-identical BOPTEST trajectories for a fixed policy, this variance is attributed to PPO/ERAM training stochasticity, not simulator noise. The single-seed canonical ($m_s\approx0.10$) is therefore the \emph{best} of five, not the median.



\paragraph{Engineering reading.} The front is strongly asymmetric. Energy-only control (0/100) collapses to a degenerate non-heating policy ($86.8\%$ violation, $m_s=1.588$): minimizing energy simply means not heating. Comfort-only (100/0) is the best single-seed $m_s$ but the most energy-hungry, and the practical 75/25 point recovers near-comfort-only comfort at lower energy --- the recommended deployment compromise. The crucial caveat is the seed dimension: the single-seed front overstates reliability, because the $N=5$ extension turns both canonicals into wide distributions (CV $0.42$--$0.61$), so the Pareto curve should be read as a best-seed envelope, not a deployment guarantee.





\subsection{Seasonal variance falsification (roadmap \S9; audit \S13)}

At $N=3$, the practical canonical appeared to show near-deterministic winter behavior and a seasonal inversion relative to the neutral canonical, motivating a pre-specified falsification test before seeds 45 and 46 were trained. The direction-specific predictions (February winter $\sigma(m_s)<0.005$; June summer $\sigma(m_s)>0.05$; winter neutral/practical variance ratio $>20$) failed at $N=5$: February practical $\sigma(m_s)$ rose to order $0.17$ and the winter variance ratio collapsed to order one. This is a success of the pre-specification protocol (audit anchors \texttt{93df9b3} pre-specification, \texttt{62dc859} post-N=5 falsification), not a project failure. The honest, narrower conclusion: MORL is promising in mean performance, especially at comfort-leaning preferences, but is \emph{not} deployment-stable without explicit stabilization (validation-based checkpoint selection, early stopping, or ensemble selection), which is left as future work because the canonical protocol fixes final-epoch evaluation.





\subsection{PI reference and synthesis (roadmap \S10--\S11)}

The BOPTEST built-in PI controller is a reproducible reference, not a tuned strong baseline. On the 12-month yearly evaluation it is comfort-poor but energy-frugal: mean RMSE$_T=3.40\,^\circ$C, mean violation $63.6\%$, mean monthly energy $104.1$ kWh, mean $m_s=0.910$. All Block 2 RL/MORL agents reach $m_s$ well below PI's $0.910$, dominating it on comfort at comparable or lower energy after accounting for the comfort/energy trade-off.

Table~\ref{tab:ms_decomp} consolidates the maintenance-score decomposition across every controller family. It makes the cross-family pattern explicit: usable controllers (pure v3, hybrid, MORL 17D) keep both $r_{\mathrm{time}}$ and $r_{\mathrm{sev}}$ small, whereas the failures (direct v3.5, PI) carry a large $r_{\mathrm{time}}$ --- they spend a large fraction of the horizon outside the band. PI is the mirror image of direct v3.5: both have $r_{\mathrm{time}}$ near or above $0.6$, but PI fails by chronic under-heating while direct v3.5 fails by saturated overshoot.

\begin{table}[pos={!ht}]
\centering
\caption{Cross-controller maintenance-score decomposition $m_s = r_{\mathrm{time}} + r_{\mathrm{sev}}$ (Eq.~\ref{eq:ms}). $r_{\mathrm{time}}$ = fraction of steps outside the band; $r_{\mathrm{sev}}$ = worst relative band exceedance. Thermostatic/HDRL rows are 14-day windows; MORL/PI rows are 12-month yearly means.}
\label{tab:ms_decomp}
\small
\begin{tabular}{lrrr}
\toprule
Controller (window/eval) & $m_s$ & $r_{\mathrm{time}}$ & $r_{\mathrm{sev}}$ \\
\midrule
pure v3 (peak) & 0.073 & 0.015 & 0.058 \\
direct v3.5 (peak) & 1.046 & 0.771 & 0.276 \\
hybrid $\lambda_T{=}0.10$ (peak) & 0.087 & 0.047 & 0.040 \\
HDRL $\lambda_T{=}0$ (peak) & 0.180 & 0.061 & 0.119 \\
pure v3 (typical) & 0.095 & 0.044 & 0.051 \\
direct v3.5 (typical) & 1.102 & 0.824 & 0.278 \\
hybrid $\lambda_T{=}0.10$ (typical) & 0.041 & 0.024 & 0.017 \\
HDRL $\lambda_T{=}0$ (typical) & 0.234 & 0.031 & 0.202 \\
MORL 50/50 (yearly, N=5) & 0.187 & 0.130 & 0.057 \\
MORL 75/25 (yearly, N=5) & 0.139 & 0.092 & 0.047 \\
PI (yearly) & 0.910 & 0.636 & 0.274 \\
\bottomrule
\end{tabular}
\end{table}

Block 2 establishes four controller-side claims. First, predictive fidelity and RL training utility are not equivalent: direct v3.5 is the more accurate twin yet fails as a rollout environment (live violation $>77\%$), and the matched-resolution ablation (Table~\ref{tab:coarse_graining}) shows the \emph{same black-box} v3 also collapses once its temporal resolution---and predictive accuracy---is increased, isolating the fidelity/smoothing trade-off rather than the model class. Second, role separation works --- v3 provides smooth rollout dynamics while frozen v3.5 acts as a physical censor through disagreement shaping --- giving the canonical hybrid ($m_s=0.087$ peak, $0.041$ typical). Third, the censor strength is controller-family specific: HDRL rejects $\lambda_T=0.10$ and is best at $\lambda_T=0$. Fourth, MORL is viable only with the 17D interface, and its N=5 analysis reveals high seed variance and falsifies the N=3 seasonal-inversion mechanism. The engineering implication is that hybrid surrogate RL is a role-allocation problem --- rollout smoothness, physical censoring, observation geometry, controller family, and seed stabilization are separate design axes. This is the bridge to Block 3, where the fixed \texttt{bestest\_air} recipe is transferred to the hydronic BOPTEST family.
